\documentclass[10pt,a4paper]{amsart}
\usepackage[margin=19mm]{geometry}
\usepackage{amsmath,amssymb,mathtools}
\usepackage[scr=rsfs,cal=euler]{mathalfa}
\usepackage{xfrac}
\usepackage[unicode,hidelinks,bookmarks=false]{hyperref}
\newcommand{\ie}{i.e.\ }
\newcommand{\cf}{cf.\ }
\newcommand{\resp}{respectively\ }
\newcommand{\Subsection}{\subsection}
\providecommand{\nobreakdash}{\nobreak\hbox{-}}
\setlength{\emergencystretch}{2em}
\multlinegap0pt
\usepackage{amssymb}
\newtheorem{thm}{Theorem}
\title{Integrality of Averages of Roots of Unity and Perfect Isometries}
\author{Chatchawan Panraksa, Pornrat Ruengrot}
\date{Source: arXiv:2605.31161v1 (CC BY 4.0)}
\begin{document}
\maketitle

\noindent\emph{Source and licence.} Integrality of Averages of Roots of Unity and Perfect Isometries. Version arXiv:2605.31161v1 (CC BY 4.0), distributed by arXiv under the Creative Commons Attribution 4.0 International licence, http://creativecommons.org/licenses/by/4.0/, as recorded in the arXiv licence metadata for this exact version and shown on the abstract page https://arxiv.org/abs/2605.31161v1. Full licence notice: \texttt{licenses/arxiv-2605.31161-CC-BY-4.0.txt}.\par\smallskip

\noindent\textit{Editorial scope.} Counted unit: the article's thm thm:perfect (source0.tex lines 255--262). The article's complete original proof of it is delivered with it (source0.tex lines 264--300, one \texttt{proof} environment of 37 lines). No mathematics is written, repaired or re-derived: the statement and the proof are the article's own lines, in the article's own notation, and no retained line is substituted or re-flowed.  Retained uncounted as the source-local context the proof needs: lines 77--85; lines 145--153.  Nothing was omitted from any retained span, so the package carries no omission record.  No retained line contains a citation, so the wrapper carries no bibliography and no bibliography entry.  No retained line contains a figure environment, an image-inclusion command or drawing code, so no image is delivered and the package declares no asset dependency.  Editorial formatting only, in addition: an AMS \texttt{amsart} preamble that restates the article's own declarations and macros used by the retained lines, namely the environments \texttt{thm} on separate counters, together with the standard packages those lines need.  The retained environments print in this document as Theorem 1 in order of appearance, whereas the article prints the same items with the numbering of its own full document.  The complete native source of the article as distributed in the arXiv e-print for this exact version is bundled unchanged as \texttt{sources/201710/source0.tex}.
\begin{thm}\label{thm:alln}
	Let $n\ge 1$, $\omega=e^{2\pi i/n}$, and $f:\mathbb Z_n\to\mathbb Z_n$.
	If for every $b\in\mathbb Z_n$ the average
	\[
		\mu_{b}(f)=\frac1n\sum_{x=0}^{n-1}\omega^{\,f(x)+bx}
	\]
	is an algebraic integer, then there exist $\alpha, \beta \in \mathbb{Z}_n$ such that
	$f(x)\equiv \alpha x+\beta \pmod n$ for all $x\in\mathbb Z_n$.
\end{thm}

\begin{thm}\label{thm:local-global}
	Let $p$ be a prime and $r\ge 1$. With notation as above, let $S\in \mathcal{O}_F$ and set $\mu=S/p^r\in F$.
	Let $\mathfrak{p}=(1-\omega)$ be the unique prime of $\mathcal{O}_F$ above $(p)$. If (under the natural embedding
	$F\hookrightarrow F_{\mathfrak{p}}\cong \mathbb{Q}_p(\omega)$) we have
	\[
		\mu\in \mathcal{O}_{F_{\mathfrak{p}}}\cong \mathbb{Z}_p[\omega],
	\]
	then $\mu\in \mathcal{O}_F$. In particular, $\mu$ is an algebraic integer of $F$.
\end{thm}

\begin{thm}\label{thm:perfect}
	Let $n=p^r$.  A bijection $f : \mathbb{Z}_n \to \mathbb{Z}_n$ gives a perfect
	isometry of the group $C_n$ if and only if there exist $\alpha,\beta \in \mathbb{Z}_n$
	with $\gcd(\alpha,n)=1$ such that
	\[
		f(x) \equiv \alpha x + \beta \pmod n \quad \text{for all } x\in\mathbb{Z}_n.
	\]
\end{thm}

\begin{proof}
	$(\Rightarrow)$\ Suppose $f$ yields a perfect isometry.  By condition~(i),
	for every $a,b\in\mathbb{Z}_n$,
	\[
		\frac{1}{n}\sum_{j=0}^{n-1}\omega^{a f(j)+bj} \in \mathbb{Z}_p[\omega].
	\]
	Fix $a=1$.  Then for all $b$,
	\[
		\frac{1}{n}\sum_{j=0}^{n-1}\omega^{f(j)+bj} \in \mathbb{Z}_p[\omega].
	\]
	By Theorem~\ref{thm:local-global} this average is an algebraic integer, so Theorem~\ref{thm:alln} applies and gives $\alpha,\beta\in\mathbb{Z}_n$ with
	\[
		f(j)\equiv \alpha j + \beta \pmod n \quad\text{for all }j.
	\]
	Since $f$ is bijective, multiplication by $\alpha$ must be invertible modulo~$n$;
	hence $\gcd(\alpha,n)=1$.

	$(\Leftarrow)$\ Conversely, assume $f(x)\equiv \alpha x+\beta$ with
	$\gcd(\alpha,n)=1$.  Then
	\[
		\mu_I(g^a,g^b)
		= \sum_{j=0}^{n-1}\omega^{a(\alpha j+\beta)+bj}
		= \omega^{a\beta}\sum_{j=0}^{n-1}\omega^{(a\alpha+b)j}.
	\]
	The inner sum equals $n$ if $a\alpha+b\equiv0\pmod n$ and $0$ otherwise, so
	\[
		\frac{\mu_I(g^a,g^b)}{n}
		= \begin{cases}
			\omega^{a\beta}\in \mathbb{Z}_p[\omega], & b\equiv -a\alpha \pmod n, \\[4pt]
			0,                                       & \text{otherwise}.
		\end{cases}
	\]
	Thus condition~(i) holds.  For~(ii), if $\mu_I(g^a,g^b)\ne0$ then
	$b\equiv -a\alpha\pmod n$, and because $\alpha$ is a unit, we have
	$a=0$ if and only if $b=0$.  Hence $g^a$ and $g^b$ are simultaneously
	the identity or both non-identity elements, satisfying the separation property.
\end{proof}

\end{document}
